% Part: proof-theory
% Chapter: proof-search
% Section: search-algorithm

\documentclass[../../../include/open-logic-section]{subfiles}

\begin{document}

\olfileid[tr]{pt}{ps}{sal}

\olsection{\Log{G3c} İçin Arama Algoritması}

\Log{G3c} için bir arama algoritması betimliyoruz. Bu, olanaklı tek
algoritma olmadığı gibi en verimlisi de değildir. Ama doğrudur: Bir
$\Gamma \Sequent \Delta$ sekantı için !!a{proof} varsa sonunda durup
!!a{proof} bulur. Ayrıca şaşmazdır: !!a{proof} bulmadan durur ya da hiç
durmazsa $\Gamma \Sequent \Delta$ baştan geçerli değildir.

Algoritmanın zorunlu temel özelliklerinden biri, $\Gamma \Sequent \Delta$
içindeki ya da ispat araması sırasında üretilen bir sekanttaki her
!!{formula} için gerektiği kadar sık ``indirgeme'' yapılmasını güvence altına
almasıdır; yani bu formül, geriye doğru uygulanan bir çıkarım kuralının
asıl formülü olarak ele alınır. Bu özelliğe \emph{adillik} deriz.

Algoritma aşamalar hâlinde çalışır. Her aşamanın çıktısı bir sekant ağacıdır;
sonraki aşamanın çıktısı ise daha önce aksiyom olmayan her en üst sekantta
!!a{formula} indirgenerek üretilir.

$0$. aşamada $\Gamma \Sequent \Delta$ ile başlarız. Adilliği güvence altına
almak üzere $\Gamma$ ve $\Delta$ içindeki !!{formula} biçimlerine, farklı
!!{formula} biçimleri farklı sayılar alacak biçimde indisler atarız.
Örneğin !!{formula} biçimlerini soldan sağa doğru saydığınızı düşünün. Sayıları üst indis
olarak yazarız. Örneğin $!A, !B \Sequent !C$ için !!a{proof} arıyorsak
$0$.~aşamanın çıktısı $!A^1, !B^2 \Sequent !C^3$ olur. Böyle indislenmiş bir
sekantta üst indis olarak geçen en büyük~$n$, sekantın \emph{en büyük indisi}dir
(örnekte~$3$).

Sekant niceleyiciler içeriyorsa sağdaki $\lexists[x][!A(x)]$ ya da soldaki
$\lforall[x][!A(x)]$ indirgenirken hangi terimlerin kullanılacağını da
bilmemiz gerekir. Yalnızca !!{constant} simgeleri $\Obj c_0$, $\Obj c_1$, $\Obj c_2$,
\dots{} ile~$\Gamma$ ve $\Delta$ içinde geçen fonksiyon simgelerinden
kurulan terimlere ihtiyaç duyarız. Bütün bu terimler kümesinin sabit bir
sıralaması verildiğini varsayarız; böylece örneğin ``$\Gamma \Sequent
\Delta$ içinde geçmeyen ilk !!{constant} simgesinden söz edebiliriz.'' Algoritmanın
ürettiği ağaçtaki her indisli sekantla ilişkilendirilmiş bir !!{constant}
simgesi kümesi vardır. $0$.~aşamadaki sekanta atanan !!{constant} simgesi kümesi, sekantta
geçen bütün !!{constant} simgelerinin kümesidir; sekant hiç !!{constant} içermiyorsa bu
küme $\{\Obj c_0\}$'dır.

Algoritmanın $n$.~aşamada, ilişkilendirilmiş !!{constant} simgesi kümeleriyle
birlikte bir indisli sekantlar ağacı ürettiğini varsayın. $n+1$.~aşamada her
en üst sekant için aşağıdakileri yaparız.

Sekantın hem solunda hem sağında geçen bir !!a{formula}~$!A$ varsa ya da
sol taraf~$\lfalse$ içeriyorsa hiçbir şey yapmayız: Bu sekantların her biri
için \Log{G3c} içinde bir !!{proof} vardır ve bu belirli en üst sekant için yapılacak
başka bir şey kalmamıştır.

Sekantta atomik olmayan hiçbir !!{formula} yoksa algoritmayı başarısız olarak
durdurun. $\Gamma \Sequent \Delta$ için ispat yoktur.

Aksi hâlde en küçük~$i$ indisine sahip atomik olmayan, $!A$ adlı !!{formula} biçimini
seçin. Söz konusu en üst sekant ya $!A^i, \Pi \Sequent \Lambda$ ya da
$\Pi \Sequent \Lambda, !A^i$ biçimindedir. $k$, sekantın en büyük indisi;
$C$ ise sekanta atanmış !!{constant} simgesi kümesi olsun.

$!A^i$'yi ``indirgeriz''; yani~$!A$ içindeki !!{main operator} ile $!A^i$'nin
solda mı sağda mı geçtiğine göre belirlenen çıkarım kuralı uyarınca yeni en
üst sekantlar üretiriz.

\begin{enumerate}
  \item $!A \ident !B \land !C$ ve solda geçiyor: $(!B \land !C)^i,
  \Pi \Sequent \Lambda$ üzerine yeni bir en üst sekant ekleyin:
  \[
  \Axiom$!B^{k+1}, !C^{k+2}, \Pi \fCenter \Lambda$
  \RightLabel{\LeftR{\land}}
  \UnaryInf$(!B \land !C)^i, \Pi \fCenter \Lambda$
  \DisplayProof
  \]
  Yeni sekanta atanan !!{constant} simgesi kümesi~$C$'dir.
  \item $!A \ident !B \land !C$ ve sağda geçiyor: $\Pi \Sequent
  \Lambda, (!B \land !C)^i$ üzerine iki yeni en üst sekant ekleyin:
  \[
  \Axiom$\Pi \fCenter \Lambda, !B^{k+1}$
  \Axiom$\Pi \fCenter \Lambda, !C^{k+1}$
  \RightLabel{\RightR{\land}}
  \BinaryInf$\Pi \fCenter \Lambda, (!B \land !C)^i$
  \DisplayProof
  \]
  Yeni sekanta atanan !!{constant} simgesi kümesi~$C$'dir.

  \item $!A \ident \lnot !B$, $!A \ident !B \lor !C$ ve $!A \ident !B
  \lif !C$ durumları benzerdir ve alıştırma olarak bırakılmıştır.

  \item $!A \ident \lexists[x][!B(x)]$ ve solda geçiyor:
  $\lexists[x][!B(x)]^i, \Pi \Sequent \Lambda$ üzerine yeni bir en üst
  sekant ekleyin:
  \[
  \Axiom$!B(c)^{k+1}, \Pi \fCenter \Lambda$
  \RightLabel{\LeftR{\lexists}}
  \UnaryInf$\lexists[x][!B(x)]^i, \Pi \fCenter \Lambda$
  \DisplayProof
  \]
  Burada $c$,~$C$ içinde bulunmayan ilk !!{constant} simgesidir. Yeni sekanta
  atanan !!{constant} simgesi kümesi $C \cup \{c\}$'dir.
  \item $!A \ident \lexists[x][!B(x)]$ ve sağda geçiyor: $\Pi \Sequent
  \Lambda, \lexists[x][!B(x)]$ üzerine yeni bir en üst sekant ekleyin:
  \[
  \Axiom$\Pi \fCenter \Lambda, \lexists[x][!B(x)]^{k+1}, !B(t)^{k+2}$
  \RightLabel{\RightR{\lexists}}
  \UnaryInf$\Pi \fCenter \Lambda, \lexists[x][!B(x)]^i$
  \DisplayProof
  \]
  Burada $t$, yalnızca~$C$ içindeki !!{constant} simgelerini içeren ve bu
  sekantın altında $\lexists[x][!B(x)]$ sağda indirgenirken daha önce
  kullanılmamış ilk terimdir (böyle terimlerin tümü kullanılmışsa $\Obj
  c_0$'dır). Yeni sekanta atanan !!{constant} simgesi kümesi~$C$'dir.
  \item $!A \ident \lforall[x][!B(x)]$ durumları benzerdir ve alıştırma
  olarak bırakılmıştır.
\end{enumerate}

$n+1$.~aşamada hiçbir en üst sekanta yeni bir sekant eklemediysek algoritma
başarıyla sona erer. Bu durumda $n+1$.~aşamadaki ağaçtan bütün indisleri
silerek elde edilen,~$\Gamma \Sequent \Delta$ için !!a{proof} bulmuşuzdur.
En üst sekantların tümü $!A, \Pi \Sequent \Lambda, !A$ ya da $\lfalse,
\Pi \Sequent \Lambda$ biçimindedir. Bütün çıkarımlar doğru \Log{G3c}
çıkarımlarıdır. Önerme mantığı çıkarımları için bu açıktır. Her adımda bir
sekanta atanan !!{constant} simgesi kümesi~$C$'nin, sekantta geçen bütün
!!{constant} simgelerini içerdiğine dikkat edin. Dolayısıyla
\LeftR{\lexists} ve \RightR{\lforall} kullanılan indirgemelerde özdeğişken
koşulu sağlanır: Özdeğişken~$c$, sonuç sekantına atanan~$C$ kümesinde
bulunmayan !!a{constant} olduğundan sonuçta da geçmez. Şunu elde ederiz:

\begin{prop}
\Log{G3c} için ispat arama algoritması doğrudur: Başarıyla sona ererse
başlangıç sekantı~$\Gamma \Sequent \Delta$ için !!a{proof} vardır.
\end{prop}






\end{document}
